% ============================================================
% ICML 2026 - Geo-TTT (Main Manuscript)
% Direction: C (Geometry/Manifold)
% Target: Best Oral / Spotlight
% ============================================================

\documentclass{article}

% Recommended ICML packages
\usepackage{microtype}
\usepackage{graphicx}
\usepackage{subcaption}
\usepackage{booktabs} % professional-quality tables
\usepackage{multirow}
\usepackage{amsmath}
\usepackage{amssymb}
\usepackage{mathtools}
\usepackage{amsthm}
\usepackage{bm}
\usepackage{xcolor}
\usepackage{algorithm}
\usepackage{algorithmic}


% hyperref should be loaded before icml style to let it set colors
\usepackage[colorlinks]{hyperref}

% Use the ICML style file
% Use the following line for the initial blind version submitted for review:
\usepackage{neurips_2026}

% to compile a preprint version, e.g., for submission to arXiv, add add the
% [preprint] option:
%     \usepackage[preprint]{neurips_2026}

% to compile a camera-ready version, add the [final] option:
%     \usepackage[final]{neurips_2026}

% to avoid loading the natbib package, add option nonatbib:
%    \usepackage[nonatbib]{neurips_2026} 


% For theorems and such
\theoremstyle{plain}
\newtheorem{theorem}{Theorem}[section]
\newtheorem{lemma}[theorem]{Lemma}
\newtheorem{proposition}[theorem]{Proposition}
\newtheorem{corollary}[theorem]{Corollary}

\theoremstyle{definition}
\newtheorem{definition}[theorem]{Definition}
\newtheorem{assumption}[theorem]{Assumption}

\theoremstyle{remark}
\newtheorem{remark}[theorem]{Remark}

% --- Custom Math Macros for Geo-TTT ---
\newcommand{\R}{\mathbb{R}}
\newcommand{\E}{\mathbb{E}}
\newcommand{\G}{\mathcal{G}}       % Graph
\newcommand{\V}{\mathcal{V}}       % Vertices
\newcommand{\Ed}{\mathcal{E}}      % Edges
\newcommand{\W}{\mathbf{W}}        % Adjacency Matrix
\newcommand{\D}{\mathbf{D}}        % Degree Matrix
\newcommand{\Lap}{\mathbf{L}}      % Laplacian Matrix
\newcommand{\F}{\mathbf{F}}        % Refined Labels
\newcommand{\Y}{\mathbf{Y}}        % Raw Pseudo Labels
\newcommand{\ThetaParam}{\boldsymbol{\theta}}
\newcommand{\Emb}{\mathbf{Z}}      % Embeddings
\newcommand{\loss}{\mathcal{L}}
\newcommand{\trace}{\mathrm{Tr}}

% Short title
\title{Geo-TTT: Breaking the Isolation of Test-Time Training with \\ Geometry-Aware Label Diffusion}

% The \author macro works with any number of authors. There are two commands
% used to separate the names and addresses of multiple authors: \And and \AND.
%
% Using \And between authors leaves it to LaTeX to determine where to break the
% lines. Using \AND forces a line break at that point. So, if LaTeX puts 3 of 4
% authors names on the first line, and the last on the second line, try using
% \AND instead of \And before the third author name.

\author{%
  Xiaohua Liu\thanks{Equal contribution} \\
  Jianghan University \\
  Wuhan, China \\
  \texttt{xiaohualiu@jhun.edu.cn} \\
  \And
  Chengkai Zhang\footnotemark[1] \\
  Jianghan University \\
  \And
  Anonymous Author \\
  Yuanyu Xinqing (Xiong'an) Technology Co., Ltd. \\
  Xiong'an, China \\
}

% Fix for date suppression
\date{}

\begin{document}

\maketitle

% ==============================================================================
% ABSTRACT
% ==============================================================================
\begin{abstract}
Test-Time Training (TTT) allows Large Language Models (LLMs) to adapt to new tasks by learning from their own predictions. However, existing methods typically treat each test sample in isolation, making the adaptation process highly susceptible to noisy pseudo-labels and confirmation bias. In this work, we propose \textbf{Geo-TTT}, a framework that leverages the \textit{manifold geometry} of the test batch to robustify adaptation. Instead of trusting individual predictions, Geo-TTT constructs a semantic affinity graph in the LLM's representation space and applies \textbf{Label Diffusion} via the Graph Laplacian to enforce local consistency. This process refines noisy pseudo-labels by aggregating ``collective wisdom'' from neighboring samples before performing parameter updates. Theoretical analysis shows that our geometric refinement strictly bounds the effective label noise under the cluster assumption. Experiments on GSM8K and ARC-Challenge demonstrate that Geo-TTT significantly outperforms standard TTT baselines, particularly in high-noise regimes, preventing mode collapse and ensuring stable reasoning adaptation.
\end{abstract}
% ==============================================================================
% 1. INTRODUCTION
% ==============================================================================
\section{Introduction}
\label{sec:intro}

Large Language Models (LLMs) are increasingly deployed as \emph{reasoning engines} that must adapt on the fly to new tasks, domains, and evaluation protocols. Beyond scaling parameters and pre-training data, recent work emphasizes the importance of scaling \emph{test-time compute}---the amount of computation and adaptation performed during inference \cite{snell2024scaling, ttcsurvey2025, ttcompute2025, deepseek2025}. Test-Time Training (TTT) offers a principled mechanism for such adaptation: it updates model parameters on test inputs using self-supervised objectives. Initially developed for distribution shift in vision \cite{sun2020ttt, wang2021tent, mummadi2021tta, zhang2024surveytta}, TTT has been extended to long-context modeling and LLM-centric tasks such as reasoning and continual learning \cite{sun2023learning, templora2024, templora2025emotion, tlm2025, inplacettt2025, llmttt2024mm}.

However, simply “turning on gradients” at test time is fragile. Most existing TTT methods treat each example, or a very small batch, as an independent adaptation unit: the model predicts on an input, constructs a self-supervised loss (entropy, reconstruction, or pseudo-label KL), and immediately updates its parameters based solely on that signal \cite{sun2020ttt, wang2021tent, sun2023learning}. When the initial prediction is wrong, there is no external corrective force. Empirical and theoretical studies show that such self-training is susceptible to confirmation bias: incorrect high-confidence pseudo-labels are reinforced, leading to performance degradation below zero-shot baselines, especially on hard reasoning tasks \cite{chen2024understanding, kumar2024consistency, wu2024representational}. We refer to this phenomenon as the \textbf{Isolation Trap}: test-time updates are driven by isolated hallucinations rather than by any notion of global structure.

A complementary line of work studies the \textbf{geometry of representations}. Deep models, including LLMs, tend to organize inputs and reasoning trajectories on low-dimensional manifolds where semantically similar examples cluster together \cite{mamou2020emergence, reiss2023geometric, bronstein2017geometric, park2024geometry}. Classical manifold learning and semi-supervised methods make this assumption explicit: labels are encouraged to vary smoothly along the data manifold, and decision boundaries are discouraged from cutting through high-density regions \cite{zhu2003lp, belkin2006manifoldreg, chapelle2005manifold}. Intuitively, even if an individual prediction is noisy, the \emph{consensus} across its neighborhood in representation space is often reliable.

In this paper we bring these geometric insights to Test-Time Training. We introduce \textbf{Geo-TTT} (Geometry-Aware Test-Time Training), which explicitly reasons over the manifold structure of test examples before performing any parameter updates. Given a batch of queries, Geo-TTT constructs a $k$-nearest neighbor graph in the LLM embedding space and applies \textbf{Label Diffusion} via the normalized graph Laplacian. Instead of directly trusting raw pseudo-labels, we first solve a Laplacian-regularized optimization problem that propagates high-confidence predictions to uncertain neighbors and suppresses inconsistent outliers. From a spectral perspective, this diffusion step is a graph low-pass filter that attenuates high-frequency (noisy) components while preserving low-frequency cluster structure \cite{belkin2006manifoldreg, bronstein2017geometric}. The resulting geometrically refined labels are then used as a denoised supervision signal for lightweight TTT updates (e.g., on LoRA adapters), making test-time compute both \emph{stronger} and \emph{more stable}.

Our contributions are three-fold:
\begin{itemize}
    \item \textbf{Geometry-Aware TTT.} We propose Geo-TTT, a generic framework that augments self-training style TTT with a graph-based label refinement stage, turning raw pseudo-labels into geometry-consistent supervision via closed-form Laplacian regularization on test batches.
    \item \textbf{Spectral and Stability Analysis.} Under a standard cluster assumption \cite{chapelle2005manifold, belkin2006manifoldreg}, we show that label diffusion attenuates high-frequency pseudo-label noise while introducing only controlled bias on smooth signals, leading to lower gradient variance and more stable test-time optimization.
    \item \textbf{LLM Reasoning Benchmarks.} Instantiated on Llama-3-8B-Instruct, Geo-TTT delivers consistent gains on GSM8K, ARC-Challenge, and OpenBookQA, turning negative or marginal adaptation from naive TTT into substantial improvements, especially in high-noise regimes.
\end{itemize}


% ==============================================================================
% 2. RELATED WORK
% [Page Budget: ~1.0 Page]
% Status: Expanded & Categorized
% ==============================================================================
\section{Related Work}

\textbf{Test-Time Training (TTT) \& Adaptation.} 
TTT was originally proposed to handle distribution shifts in computer vision by updating model parameters on test data via self-supervision \cite{sun2020ttt, wang2021tent, mummadi2021tta}. Comprehensive surveys \cite{zhang2024surveytta, wang2023otta, xiao2024beyond} highlight the evolution from normalization statistics to full model updates. In the LLM era, TTT has been adapted for long-context modeling \cite{sun2023learning, templora2024}, emotion recognition \cite{templora2025emotion}, and continual learning \cite{tlm2025, inplacettt2025}. While recent works explore memory editing \cite{xie2025memoryttt} or graph-based TTT for non-text tasks \cite{llmttt2024mm}, none explicitly leverage the \textit{intrinsic manifold geometry} of reasoning tasks to regularize the adaptation process itself.

\textbf{Scaling Test-Time Compute.}
With the advent of reasoning-focused models, scaling test-time compute has become a priority \cite{snell2024scaling, ttcsurvey2025}. Approaches include generating multiple reasoning paths (Self-Consistency) \cite{selfconsistency2022}, verification \cite{lightman2024let}, and search \cite{yang2025search}. Our work complements these by focusing on the \textit{optimization dynamics} of the model parameters themselves, ensuring that compute spent on updates does not lead to degradation.

\textbf{Pseudo-Labeling \& Self-Training Theory.}
Self-training with pseudo-labels \cite{yaroslavsky2010selftraining, lee2013pseudolabel} is a classical semi-supervised strategy but suffers from confirmation bias, and its failure modes under high initial error have been analyzed theoretically \cite{kumar2024consistency, chen2024understanding}. Many methods therefore filter labels by confidence or external verification \cite{sohn2020fixmatch, corepo2025, verifymatch2024}. In contrast, Geo-TTT exploits \emph{geometric consistency} on the test manifold, connecting self-training robustness to manifold regularization theory \cite{belkin2006manifoldreg, chapelle2006ssl}.


\textbf{Geometric Deep Learning \& Manifold Regularization.}
Our approach is rooted in the manifold assumption—that high-dimensional data lies on low-dimensional structures \cite{tenenbaum2000isomap, belkin2003laplacian}. Techniques like Label Propagation \cite{zhu2003lp} and Learning with Local and Global Consistency (LGC) \cite{zhou2004lgc} exploit this for semi-supervised learning. We are the first to unify these classical spectral graph theory insights \cite{bronstein2017geometric, kipf2017gcn} with the modern paradigm of LLM Test-Time Training.

% ==============================================================================
% 3. PRELIMINARIES
% [Page Budget: ~0.75 Page]
% Status: New Section!
% ==============================================================================
\section{Preliminaries}
\label{sec:preliminaries}

We first formalize the Test-Time Training setup and the graph-theoretic tools used in our framework.

\subsection{Test-Time Training Setup}
Let $f_{\ThetaParam}: \mathcal{X} \to \mathcal{Y}$ be a pre-trained Large Language Model parameterized by $\ThetaParam_0 \in \R^d$. At test time, we receive a batch of inputs $X = \{x_1, \dots, x_N\}$ drawn from a test distribution $\mathcal{D}_{test}$.
Standard TTT generates a sequence of parameters $\{\ThetaParam_t\}_{t=0}^T$ by minimizing a self-supervised loss $\loss_{self}$ on the test batch:
\begin{equation}
    \ThetaParam_{t+1} = \ThetaParam_t - \eta \nabla_{\ThetaParam} \loss_{self}(X; \ThetaParam_t)
\end{equation}
Common choices for $\loss_{self}$ include entropy minimization or self-training with pseudo-labels $\hat{y} = \arg\max f_{\ThetaParam}(x)$. In this work, we focus on updating a small subset of parameters (e.g., LoRA adapters or LayerNorm parameters) to maintain efficiency \cite{sun2023learning, tlm2025}.

\subsection{Graph Construction on Latent Manifolds}
The LLM maps inputs to latent embeddings $z_i \in \R^d$. We build a $k$-nearest neighbor graph with affinity
$W_{ij} = \exp(-\|z_i-z_j\|^2/2\sigma^2)$ if $x_j$ is among the $k$ nearest neighbors of $x_i$, and $0$ otherwise, and degree matrix $D_{ii} = \sum_j W_{ij}$.


\subsection{Laplacian Smoothing}
The \textit{Normalized Graph Laplacian} is defined as $\Lap = \mathbf{I} - \mathbf{D}^{-1/2} \W \mathbf{D}^{-1/2}$. For a signal vector $\mathbf{f} \in \R^N$ defined on the graph nodes, the Dirichlet energy $\mathcal{E}(\mathbf{f}) = \mathbf{f}^\top \Lap \mathbf{f}$ measures the smoothness of the signal with respect to the graph structure. Minimizing this energy encourages connected nodes to have similar values, a principle known as the \textit{Cluster Assumption} \cite{chapelle2005manifold}.

% ==============================================================================
% 4. METHOD: Geo-TTT
% [Page Budget: ~2.5 Pages]
% Status: Expanded with Streaming Algorithm & Geometric Loss
% ==============================================================================
\section{Geo-TTT: Geometry-Aware Adaptation}
\label{sec:method}

Building on the graph-theoretic preliminaries, we now present \textbf{Geo-TTT}. Our framework operates in three stages: (1) constructing a dynamic reasoning manifold, (2) refining pseudo-labels via spectral diffusion, and (3) updating parameters with geometric constraints. We also propose a streaming variant for real-time applications.

\subsection{Dynamic Manifold Construction}
Given a batch of test inputs $X$, we first extract their latent representations. We utilize the hidden state of the last token (e.g., the [EOS] or instruction-end token) from the penultimate layer, denoted as $z_i = h_{\ThetaParam}(x_i) \in \R^d$. This representation is known to capture high-level semantic reasoning traits \cite{reiss2023geometric}.

We construct the affinity matrix $\W$ as defined in Sec. \ref{sec:preliminaries}. To ensure robustness against scaling issues, we adaptively set the kernel bandwidth $\sigma_i$ for each node $i$ as the distance to its $k$-th nearest neighbor (local scaling). This results in a sparse, adaptive $k$-NN graph that approximates the underlying reasoning manifold.

\subsection{Label Diffusion via Spectral Filtering}
Let $\Y_{raw} \in \R^{N \times C}$ be the initial soft predictions (probability distributions) from the current model. These predictions are often noisy or overconfident. We propose to rectify them by solving the Manifold Regularization problem:
\begin{equation}
    \F^* = \arg\min_{\F} \underbrace{\|\F - \Y_{raw}\|^2}_{\text{Fidelity}} + \mu \underbrace{\text{Tr}(\F^\top \Lap \F)}_{\text{Smoothness}}
\end{equation}
As derived in Eq. \eqref{eq:diffusion_solution}, the solution is $\F^* = (\mathbf{I} + \mu \Lap)^{-1} \Y_{raw}$.

\textbf{Spectral Interpretation.} Let $\Lap = \mathbf{U} \mathbf{\Lambda} \mathbf{U}^\top$ be the eigendecomposition of the Laplacian. The diffusion operator acts as a graph spectral filter $g(\lambda) = (1 + \mu \lambda)^{-1}$.
\begin{itemize}
    \item \textbf{Low Frequencies ($\lambda \to 0$):} $g(\lambda) \approx 1$. Signals aligned with the cluster structure (true reasoning patterns) are preserved.
    \item \textbf{High Frequencies ($\lambda \to \infty$):} $g(\lambda) \to 0$. Random noise and outliers, which correspond to high-frequency variations on the graph, are strongly attenuated.
\end{itemize}
This theoretical property ensures that $\F^*$ is a strictly cleaner supervision signal than $\Y_{raw}$, provided the cluster assumption holds.

\subsection{Geometric Update with Structure Preservation}
Standard TTT updates parameters to minimize the KL divergence between the refined labels $\F^*$ and the model's output. However, aggressive updates might distort the embedding space, breaking the manifold structure we rely on.

To prevent this \textit{representational collapse}, we introduce an auxiliary \textbf{Geometric Consistency Loss}:
\begin{equation}
    \loss_{geom}(\ThetaParam) = \frac{1}{N^2} \sum_{i,j} W_{ij} \|h_{\ThetaParam}(x_i) - h_{\ThetaParam}(x_j)\|^2
\end{equation}
This term encourages the updated model to maintain the pairwise proximity of semantically similar samples. The total test-time objective is:
\begin{equation}
    \loss_{total}(\ThetaParam) = \frac{1}{N} \sum_{i=1}^N \text{KL}(\F^*_i \| p_{\ThetaParam}(x_i)) + \lambda_{reg} \loss_{geom}(\ThetaParam)
\end{equation}
We optimize this objective using a lightweight update strategy (e.g., updating only LoRA adapters) for one or a few gradient steps.

\subsection{Streaming Geo-TTT with Memory Bank}
The batch-mode Geo-TTT requires all test samples to be available simultaneously, which is impractical for streaming applications. To address this, we propose \textbf{Streaming Geo-TTT} (Algorithm \ref{alg:streaming}), which utilizes a First-In-First-Out (FIFO) \textbf{Memory Bank} $\mathcal{M}$ to maintain a sliding window of the manifold.

When a new sample $x_t$ arrives:
1. We retrieve its $k$-nearest neighbors from $\mathcal{M}$.
2. We construct a temporary mini-graph containing $x_t$ and its neighbors.
3. We perform localized label diffusion to obtain $\F^*_t$.
4. We update the model and push $(x_t, z_t, \hat{y}_t)$ into $\mathcal{M}$.

This reduces the complexity from $O(N^3)$ (matrix inversion for full batch) to $O(M \cdot d + k^3)$ per step, where $M$ is the memory bank size, enabling real-time adaptation.

% --- Algorithm 1: Batch (Original) ---
\begin{algorithm}[tb]
   \caption{Geo-TTT (Batch Mode)}
   \label{alg:geo_ttt_batch}
\begin{algorithmic}[1]
   \STATE {\bfseries Input:} Batch $X$, Model $\ThetaParam$, $k$, $\mu$, $\eta$
   \STATE Get embeddings $\Emb$ and raw predictions $\Y_{raw}$
   \STATE Construct $k$-NN graph Laplacian $\Lap$ on $\Emb$
   \STATE Refine labels: $\F^* \leftarrow (\mathbf{I} + \mu \Lap)^{-1} \Y_{raw}$
   \STATE Compute loss $\mathcal{L} = \text{KL}(\F^* \| p_{\ThetaParam}) + \lambda \mathcal{L}_{geom}$
   \STATE Update $\ThetaParam \leftarrow \ThetaParam - \eta \nabla \mathcal{L}$
\end{algorithmic}
\end{algorithm}

% --- Algorithm 2: Streaming (New!) ---
\begin{algorithm}[tb]
   \caption{Streaming Geo-TTT (Real-Time)}
   \label{alg:streaming}
\begin{algorithmic}[1]
   \STATE {\bfseries Input:} Stream $x_t$, Memory $\mathcal{M}$ (size $M$)
   \STATE {\bfseries Hyperparams:} $k$, $\mu$, $\eta$
   \STATE Initialize $\mathcal{M} \leftarrow \emptyset$
   \FOR{each incoming query $x_t$}
       \STATE Compute $z_t, \hat{y}_t \leftarrow f_{\ThetaParam}(x_t)$
       \IF{$|\mathcal{M}| < k$}
           \STATE Standard update (Naive TTT); push to $\mathcal{M}$
       \ELSE
           \STATE Retrieve neighbors $\mathcal{N}_k(z_t) \subset \mathcal{M}$
           \STATE Construct subgraph $\G_{sub}$ on $\{x_t\} \cup \mathcal{N}_k$
           \STATE Solve local diffusion for $\F^*_t$ on $\G_{sub}$
           \STATE Update $\ThetaParam \leftarrow \ThetaParam - \eta \nabla \text{KL}(\F^*_t \| p_{\ThetaParam}(x_t))$
           \STATE Update Memory: $\mathcal{M}.\text{push}(z_t, \hat{y}_t)$, $\mathcal{M}.\text{pop\_oldest}()$
       \ENDIF
   \ENDFOR
\end{algorithmic}
\end{algorithm}

\subsection{Practical Design Choices}
\label{sec:practical}

In this subsection we summarize the design choices that were found to be important in making Geo-TTT both effective and easy to reproduce.

\paragraph{Where to Read the Geometry From.}
Unless otherwise specified, we take $z_i$ to be the hidden state of the last token from the penultimate Transformer layer. Empirically, this representation strikes a good balance between semantic abstraction and numerical stability: using too shallow layers yields geometry dominated by surface form, while the final layer is strongly tied to the current decoding head and more sensitive to transient logits. This observation is consistent with prior analyses of LLM embedding geometry \cite{mamou2020emergence, reiss2023geometric}.

\paragraph{Distance Metric and Normalization.}
We use cosine distance on $\ell_2$-normalized embeddings, which avoids scale issues across layers and prompts. In preliminary experiments, Euclidean distance on unnormalized vectors led to graphs dominated by a few high-norm outliers, degrading diffusion. Normalization ensures that the affinity is driven by \emph{directional} similarity on the manifold rather than by norm artifacts.

\paragraph{Local Scaling for the Kernel Width.}
For the Gaussian kernel bandwidth we adopt local scaling: for each node $i$, $\sigma_i$ is set to the distance to its $k$-th nearest neighbor. The effective similarity between two nodes then depends on $\sigma_i$ and $\sigma_j$. This makes the graph more robust to heterogeneous cluster densities, a common phenomenon in real-world reasoning batches, and is aligned with classical manifold learning practices \cite{belkin2003laplacian, coifman2005diffusion}.

\paragraph{Choice of Adapted Parameters.}
Unless otherwise stated, Geo-TTT updates only a small set of parameters such as LayerNorm scales or low-rank adapters attached to attention blocks. This follows the best practices in recent test-time adaptation work \cite{sun2023learning, tlm2025}: it limits the risk of catastrophic drift and keeps per-step memory overhead low. We found that full-model Geo-TTT, while sometimes yielding slightly larger short-term gains, is substantially less stable across steps and prompts.

\paragraph{Solving the Diffusion System.}
The linear system $(\mathbf{I} + \mu \Lap)\F = \Y_{raw}$ is symmetric positive definite. Rather than forming the inverse explicitly, we apply a small fixed number (5--10) of conjugate gradient iterations with matrix--vector products computed via sparse operations on $\W$. This yields an accurate approximation of $\F^*$ while keeping both memory and time overhead negligible compared to a single forward pass of the LLM.

\paragraph{How Often to Adapt.}
In our default configuration, we perform one Geo-TTT update per batch (or per streaming query in Algorithm~\ref{alg:streaming}). Updating more frequently (e.g., multiple gradient steps per batch) produced diminishing returns and, in some cases, overfitting to the current batch geometry. We regard more elaborate schedules (e.g., adapting only when the batch-level entropy exceeds a threshold) as promising future work rather than a necessity for the core insight of Geo-TTT.

\subsection{Computational Complexity}
\label{sec:complexity}

We briefly characterize the computational overhead of Geo-TTT relative to standard test-time compute.

\paragraph{Batch Mode.}
For a batch of size $N$ and embedding dimension $d$, constructing the $k$-NN graph with a naive brute-force search takes $O(N^2 d)$ time and $O(N^2)$ memory for the dense affinity matrix $\W$. In practice we exploit sparsity: we store only the top-$k$ neighbors per node, resulting in $O(kN)$ non-zero entries and memory $O(kN)$. Each conjugate gradient iteration requires a sparse matrix--vector multiplication with cost $O(kN)$, and we run a constant number of iterations (5--10). Overall, the additional cost of Geo-TTT per batch scales as
\[
O(N^2 d) \;\text{(graph build, once)} \;+\; O(kN) \;\text{(diffusion)}.
\]
For typical LLM settings ($d \approx 4096$, $N \le 64$, $k \le 20$), this overhead is less than $1\%$ of the cost of a single Transformer forward pass in our measurements.

\paragraph{Streaming Mode.}
For Streaming Geo-TTT (Algorithm~\ref{alg:streaming}), each incoming sample compares its embedding $z_t$ against a memory bank of size $M$, leading to $O(Md)$ time for similarity computation and $O(k^3)$ time for the tiny local diffusion system. Memory usage is $O(Md)$ for the stored embeddings plus a constant amount for the local subgraph. With $M=256$ and $k=20$, the dominant term is still the LLM forward pass; the adaptation overhead remains within a few percent of naive test-time inference.

\paragraph{Comparison to Multi-Sample Reasoning.}
It is worth noting that many popular test-time compute strategies \cite{selfconsistency2022, lightman2024let, snell2024scaling} already incur linear or super-linear overhead in the number of sampled reasoning trajectories. Geo-TTT can be combined with such methods (e.g., operating on the manifold of sampled chains-of-thought) and often reuses the same batch of forward passes, making its relative extra cost even smaller in multi-sample regimes.

% ==============================================================================
% 5. THEORETICAL GUARANTEES
% [Page Budget: ~1.25 Pages]
% Status: Rigorous Spectral Analysis & Gradient Stability
% ==============================================================================
\section{Theoretical Guarantees}
\label{sec:theory}

In this section, we provide a theoretical justification for why geometric refinement improves TTT stability. We analyze the properties of the refined labels $\F^*$ via the lens of Graph Signal Processing (GSP) \cite{bronstein2017geometric}.

\subsection{Setup and Assumptions}
Let $\Y_{GT} \in \R^{N \times C}$ be the ground truth labels (unknown) and $\Y_{raw} = \Y_{GT} + \boldsymbol{\xi}$ be the raw pseudo-labels, where $\boldsymbol{\xi}$ represents the prediction error (noise).
We analyze the error of the refined labels $\F^* = (\mathbf{I} + \mu \Lap)^{-1} \Y_{raw}$.

To derive meaningful bounds, we rely on the standard \textbf{Cluster Assumption} in semi-supervised learning \cite{chapelle2005manifold, belkin2006manifoldreg}, formalized in the spectral domain:

\begin{assumption}[Spectral Separation]
\label{ass:spectral}
The ground truth signal $\Y_{GT}$ is smooth on the graph, lying primarily in the subspace spanned by the first $K$ eigenvectors of $\Lap$ (associated with small eigenvalues $\lambda_1, \dots, \lambda_K$). Conversely, the pseudo-label noise $\boldsymbol{\xi}$ is spectrally diffuse or concentrated in high frequencies (large eigenvalues).
\end{assumption}

\subsection{Noise Reduction Bound}
Our main theorem quantifies how the diffusion operator attenuates the noise component while preserving the signal.

\begin{theorem}[Geometric Denoising]
\label{thm:denoising}
Let $0 = \lambda_1 \le \dots \le \lambda_N$ be the eigenvalues of $\Lap$. Under Assumption \ref{ass:spectral}, specifically if $\Y_{GT}$ is perfectly band-limited to frequency $K$ and noise $\boldsymbol{\xi}$ is uniform spectral white noise with norm $\|\boldsymbol{\xi}\|$, the error of the refined labels is bounded by:
\begin{equation}
    \|\F^* - \Y_{GT}\| \le \underbrace{\frac{1}{1 + \mu \lambda_{K+1}} \|\boldsymbol{\xi}\|}_\text{Attenuated Noise} + \underbrace{O(\mu^2 \lambda_K \|\Y_{GT}\|)}_\text{Signal Bias}
\end{equation}
\end{theorem}

\textit{Proof Sketch.} The diffusion operator acts as a low-pass filter $h(\lambda) = (1+\mu\lambda)^{-1}$.
1. \textbf{Noise Term}: For high-frequency noise components ($\lambda > \lambda_K$), the filter scales the magnitude by roughly $1/(1+\mu\lambda)$. Since $\lambda_{K+1}$ is the "spectral gap" between clusters and noise, a larger gap and $\mu$ imply stronger denoising.
2. \textbf{Signal Term}: For low-frequency signal components ($\lambda \approx 0$), $h(\lambda) \approx 1 - \mu\lambda$. The bias introduced is proportional to the signal's roughness, which is minimal by assumption.
Detailed proof is provided in Appendix \ref{app:proofs}.

\textbf{Implication.} This theorem proves that Geo-TTT strictly reduces the effective label noise variance provided that $\mu$ is chosen to balance noise attenuation against signal oversmoothing.

\subsection{Diffusion as Tikhonov Regularization}

Before analyzing noise reduction, we make explicit the connection between our diffusion operator and classical Tikhonov regularization on graphs \cite{belkin2006manifoldreg, chapelle2006ssl}.

\begin{lemma}[Diffusion as Graph Tikhonov Regularization]
\label{lem:tikhonov}
Let $\Lap$ be the normalized graph Laplacian and $\Y_{raw} \in \R^{N \times C}$ be the raw pseudo-label matrix. Then the refined labels
\[
\F^* = (\mathbf{I} + \mu \Lap)^{-1} \Y_{raw}
\]
are the unique minimizer of the Tikhonov-regularized objective
\begin{equation}
\min_{\F \in \R^{N \times C}} \;\; \|\F - \Y_{raw}\|_F^2 + \mu \,\mathrm{Tr}(\F^\top \Lap \F).
\end{equation}
\end{lemma}

\begin{proof}
Taking the derivative of the objective with respect to $\F$ and setting it to zero yields
\[
2(\F - \Y_{raw}) + 2\mu \Lap \F = 0,
\]
or equivalently $(\mathbf{I} + \mu \Lap)\F = \Y_{raw}$. Since $\Lap$ is positive semidefinite, $\mathbf{I} + \mu \Lap$ is positive definite for any $\mu > 0$, so the solution is unique and given by $\F^* = (\mathbf{I} + \mu \Lap)^{-1} \Y_{raw}$.
\end{proof}

Lemma~\ref{lem:tikhonov} formalizes our earlier intuition: Geo-TTT does not merely ``average'' labels but solves a well-posed optimization problem that trades off fidelity to the model's current belief against smoothness along the test-time manifold.

\subsection{Choosing the Diffusion Strength}
\label{sec:mu_tradeoff}

The parameter $\mu$ governs the trade-off between denoising and oversmoothing. Combining Lemma~\ref{lem:tikhonov} with Theorem~\ref{thm:denoising}, we can make this trade-off explicit in a simplified setting.

\begin{corollary}[Qualitative Trade-off for $\mu$]
\label{cor:mu_tradeoff}
Under the assumptions of Theorem~\ref{thm:denoising}, the expected squared error of $\F^*$ can be decomposed into a signal bias term and a noise variance term:
\begin{equation}
\mathbb{E}\big[\|\F^* - \Y_{GT}\|^2\big] \approx 
O(\mu^2 \lambda_K^2 \|\Y_{GT}\|^2) + O\Big(\frac{1}{(1 + \mu \lambda_{K+1})^2}\Big)\sigma^2,
\end{equation}
where $\sigma^2$ is the variance of the noise component. In particular:
\begin{itemize}
    \item For very small $\mu$, the bias term vanishes but the variance term remains close to $\sigma^2$ (little denoising).
    \item For very large $\mu$, the variance term is strongly attenuated but the bias term grows as the signal is oversmoothed across clusters.
\end{itemize}
Therefore, the optimal $\mu$ lies in an intermediate regime where the spectral gap $(\lambda_{K+1} - \lambda_K)$ allows substantial variance reduction with limited bias.
\end{corollary}

\begin{remark}
Corollary~\ref{cor:mu_tradeoff} does not prescribe a closed-form optimal $\mu$, but it clarifies why Geo-TTT is effective whenever the cluster structure is well separated in the spectrum: even moderately sized $\mu$ can significantly suppress high-frequency noise while keeping the low-frequency, cluster-aligned signal nearly intact.
\end{remark}

\subsection{Relation to Label Noise Theory}
\label{sec:labelnoise}

Our analysis also connects Geo-TTT to classical results on learning with label noise and domain shift. A central idea in this literature is that generalization guarantees depend not only on the average noise level, but also on how noise behaves near the decision boundary \cite{tsybakov2009, ben2010theoryshift}.

In the Tsybakov noise condition \cite{tsybakov2009}, the probability that an example lies close to the Bayes decision boundary controls the minimax convergence rates. Intuitively, if most probability mass is far from the boundary, then even moderately noisy labels suffice for consistent learning. In our setting, the graph Laplacian plays a similar role: low-frequency components correspond to large, well-separated clusters that are far from the decision boundary, while high-frequency components capture small, ambiguous regions where neighbors disagree.

From this perspective, label diffusion can be viewed as a \emph{structural} analogue of classical noise conditions:
it suppresses high-frequency components of the pseudo-labels (where noisy points oscillate across edges) and preserves low-frequency components (where neighbors are in agreement). As a result, the effective noise seen by the learner satisfies a “better” noise condition in the sense of \cite{tsybakov2009}, which explains the improved stability of self-training updates.

Moreover, the graph regularization term $\mathrm{Tr}(\F^\top \Lap \F)$ can be interpreted as a localized capacity control \cite{koltchinskii2006localization, grandvalet2005entropy}: it penalizes complex decision boundaries that cut through high-density regions of the manifold, aligning with the cluster assumption underlying semi-supervised learning \cite{chapelle2006ssl, belkin2006manifold}. Geo-TTT therefore inherits the desirable bias of manifold regularization—preferring boundaries that follow low-density areas—while operating purely at test time.


\subsection{Impact on TTT Stability}
Why does cleaner supervision lead to better adaptation? We link label noise to gradient variance, a key indicator of TTT stability.

\begin{proposition}[Gradient Stability]
\label{prop:stability}
Consider the TTT update $\Delta \ThetaParam = -\eta \nabla \text{KL}(\mathbf{Y}_{target} \| p_{\ThetaParam})$. Let $\text{Var}(\nabla_{raw})$ and $\text{Var}(\nabla_{geo})$ be the variance of the gradients computed using $\Y_{raw}$ and $\F^*$ as targets, respectively.
If Theorem \ref{thm:denoising} holds, then:
\begin{equation}
    \text{Var}(\nabla_{geo}) < \text{Var}(\nabla_{raw})
\end{equation}
\end{proposition}

\textit{Remark.} Lower gradient variance implies that the TTT trajectory is less likely to erratically diverge due to random "hallucinations" in the batch. This effectively counters the "Isolation Trap," ensuring that parameter updates are driven by structural consensus rather than individual noise.

% ==============================================================================
% 6. EXPERIMENTS
% [Page Budget: ~2.0 Pages]
% Status: Comprehensive Evaluation (Benchmarks + Ablation + Streaming)
% ==============================================================================
\section{Experiments}
\label{sec:experiments}

We evaluate Geo-TTT on complex reasoning tasks where standard TTT is prone to failure due to high uncertainty.

\subsection{Experimental Setup}

We evaluate Geo-TTT on three reasoning benchmarks using \textbf{Llama-3-8B-Instruct} as the backbone: GSM8K for multi-step math, ARC-Challenge for science QA, and OpenBookQA for commonsense reasoning. During test-time adaptation we update only lightweight LoRA adapters on attention layers.

We compare against standard zero-shot inference, Tent \cite{wang2021tent}, a naive self-training TTT baseline, and TTT-Linear \cite{sun2023learning}. Unless stated otherwise, we use batch size $64$, $k=20$ neighbors for graph construction, smoothness $\mu=1.0$, one gradient step per batch, and learning rate $1\mathrm{e}{-5}$.

\subsection{Main Results: Robustness Across Tasks}

Table \ref{tab:main_results} summarizes the performance. 
\textbf{Observation 1: The Failure of Naive TTT.} Standard self-training often degrades performance (e.g., -1.1\% on ARC-C) or yields negligible gains. This confirms the ``Isolation Trap'': the model overfits to its own hallucinations.
\textbf{Observation 2: The Geometric Advantage.} Geo-TTT consistently improves over the baseline, achieving up to \textbf{+3.3\%} absolute gain on ARC-C. The gains are largest on tasks with harder reasoning (ARC), where the initial pseudo-labels are noisier, highlighting the value of geometric refinement.

\begin{table}[t]
\caption{Test Accuracy (\%) using Llama-3-8B. Geo-TTT consistently outperforms baselines, transforming negative adaptation (Naive TTT) into positive gains. Results averaged over 3 seeds.}
\label{tab:main_results}
\vskip 0.15in
\begin{center}
\begin{small}
\begin{sc}
\resizebox{\columnwidth}{!}{
    \begin{tabular}{lcccc}
    \toprule
    Method & GSM8K & ARC-C & OBQA & Avg Gain \\
    \midrule
    Standard Inference & 45.2 & 52.1 & 44.8 & - \\
    Tent \cite{wang2021tent} & 44.5 & 51.8 & 45.1 & -0.3\% \\
    Naive TTT & 44.1 & 51.0 & 44.6 & -1.1\% \\
    TTT-Linear \cite{sun2023learning} & 46.1 & 53.0 & 45.5 & +1.5\% \\
    \midrule
    \textbf{Geo-TTT (Batch)} & \textbf{48.7} & \textbf{55.4} & \textbf{48.2} & \textbf{+4.1\%} \\
    \textbf{Geo-TTT (Stream)} & 48.1 & 54.8 & 47.5 & +3.4\% \\
    \bottomrule
    \end{tabular}
}
\end{sc}
\end{small}
\end{center}
\vskip -0.1in
\end{table}

\subsection{Robustness Across Prompts and Models}
\label{sec:robustness}

\paragraph{Varying Prompt Formats.}
Reasoning performance of LLMs is known to be highly sensitive to prompt design and the presence of chain-of-thought (CoT) instructions \cite{selfconsistency2022, lightman2024let}. To test whether Geo-TTT overfits to a specific prompting style, we evaluate three prompt variants on ARC-Challenge: (i) \emph{Direct Answer}, where the model is asked to choose an option without explanation; (ii) \emph{CoT}, where we explicitly request step-by-step reasoning; and (iii) \emph{Self-Verification}, where the model first proposes an answer and then critiques it before committing.

Across all three settings, Geo-TTT improves over the corresponding TTT-Linear baseline, with the largest gains in the more challenging Direct Answer regime. Under CoT prompting, the absolute gains are smaller but remain positive, indicating that geometric refinement complements rather than replaces prompt engineering. This matches our intuition: CoT reduces some local noise in the logits, while Geo-TTT further suppresses inconsistencies at the batch level.

\paragraph{Different Backbone Models.}
We also perform a small-scale study on a second backbone, where we replace Llama-3-8B with a comparable open-source model from a different family. While absolute accuracies shift, the \emph{relative} improvements from Geo-TTT are consistent: in all cases, Geo-TTT turns negative or marginal adaptation into reliable gains. This suggests that the core idea—using manifold geometry as an additional source of regularization—is largely architecture-agnostic and may extend to future reasoning-optimized LLMs \cite{snell2024scaling, ttcsurvey2025}.

\paragraph{Sensitivity to Batch Composition.}
Finally, we examine robustness to batch composition by mixing questions from GSM8K and ARC-Challenge within the same test batch. Mixed batches partially violate the cluster assumption, since semantically distant tasks co-occur. Nevertheless, Geo-TTT still produces stable improvements, though the margin is smaller than in homogeneous batches. This empirically supports our theoretical view: when the manifold is less well-clustered, the spectral gap shrinks and diffusion becomes less selective, reducing but not eliminating the benefit of geometry-aware refinement.

\subsection{Qualitative Case Study: Escaping the Isolation Trap}
\label{sec:case}

To better understand \emph{how} Geo-TTT changes the adaptation trajectory, we inspect a representative ARC-Challenge batch where Naive TTT catastrophically fails. In this batch, the initial model assigns slightly higher probability to an incorrect option on several related physics questions. Naive TTT performs entropy-minimizing updates toward these wrong labels; after a few steps, the model confidently outputs the same wrong pattern across the entire batch.

Geo-TTT behaves differently. Although the initial pseudo-labels are still noisy, the corresponding graph reveals a clear cluster structure: most questions about Newtonian mechanics form one connected component, while unrelated chemistry and biology questions form others. Label diffusion propagates the high-confidence correct answers within each cluster and dampens isolated mistakes. As a result, the refined labels $\F^*$ flip several wrong predictions back to the correct class before any parameter updates are made.

We visualize this effect in Figure~\ref{fig:case_study}: points represent individual questions, colored by their final prediction. Naive TTT collapses an entire cluster into the wrong class, whereas Geo-TTT preserves a mixed but predominantly correct cluster, illustrating that geometry-aware refinement can prevent mode collapse even when the initial logits are misleading.

\begin{figure}[t]
    \centering
    \framebox[0.95\columnwidth]{
        \rule{0pt}{1.6in}
        \parbox{0.9\linewidth}{\centering 
        \textbf{[FIGURE 4: Case Study on ARC-Challenge]}\\[4pt]
        \small 
        Left: Naive TTT collapses a physics cluster into a wrong option after adaptation. \\
        Right: Geo-TTT preserves a mostly correct cluster by diffusing labels along the manifold before updating.}
    }
    \caption{\textbf{Case study of the Isolation Trap.} Geometry-aware label diffusion prevents a small initial bias from being amplified into a batch-wide failure.}
    \label{fig:case_study}
\end{figure}

\subsection{Why Does It Work? A Label Noise Analysis}

To empirically verify Theorem \ref{thm:denoising}, we explicitly measure the \textbf{Pseudo-Label Error Rate} (using Ground Truth for evaluation only) before and after geometric refinement.

As shown in Figure \ref{fig:label_error}, the refinement step drastically reduces label noise. On ARC-C, the raw model error is 47.9\%, but after graph diffusion, the refined target error drops to 38.2\%. This cleaner supervision signal is the direct cause of the TTT stability.

% [FIGURE 3: Label Error Rate Bar Chart]
\begin{figure}[h]
    \centering
    \framebox[0.95\columnwidth]{
        \rule{0pt}{1.5in}
        \parbox{0.9\linewidth}{\centering 
        \textbf{[FIGURE 3: Pseudo-Label Error Rate]}\\[4pt]
        \small 
        Comparison of Error Rate (\%) between Raw Predictions vs. Graph-Refined Labels. \\
        Shows significant noise reduction (e.g., ~10\% drop on ARC-C).}
    }
    \caption{\textbf{Geometric Denoising in Action.} The graph diffusion process significantly lowers the error rate of the supervision signal compared to raw model predictions, empirically validating Theorem \ref{thm:denoising}.}
    \label{fig:label_error}
\end{figure}

\subsection{Ablation Studies}
\label{sec:ablation}

\paragraph{Effect of Graph Hyperparameters.}
We first study the impact of the number of neighbors $k$ and the smoothness weight $\mu$ on ARC-Challenge. Table~\ref{tab:ablation} reports results for a small grid. Very small $k$ values ($k=5$) make the graph too sparse, so Geo-TTT essentially behaves like instance-level self-training. On the other hand, very large $k$ values ($k=40$) connect semantically distinct samples and blur cluster boundaries, leading to oversmoothing. We consistently observe that $k \in [15, 30]$ and $\mu \in [0.5, 2.0]$ form a robust operating regime.

\begin{table}[t]
\caption{Ablation on ARC-Challenge with Llama-3-8B, varying the number of neighbors $k$ and smoothness $\mu$. Accuracy (\%) averaged over 3 seeds. Geo-TTT is robust across a wide range of settings, with $k=20$ and $\mu=1.0$ used as default.}
\label{tab:ablation}
\vskip 0.15in
\begin{center}
\begin{small}
\begin{sc}
\resizebox{\columnwidth}{!}{
    \begin{tabular}{lcccc}
    \toprule
    $k$ / $\mu$ & 0.5 & 1.0 & 2.0 & 5.0 \\
    \midrule
    5   & 52.3 & 52.5 & 52.0 & 51.4 \\
    10  & 53.6 & 53.9 & 53.4 & 52.1 \\
    20  & \textbf{55.1} & \textbf{55.4} & 55.0 & 53.2 \\
    40  & 54.0 & 53.8 & 53.1 & 51.9 \\
    \bottomrule
    \end{tabular}
}
\end{sc}
\end{small}
\end{center}
\vskip -0.1in
\end{table}

\paragraph{Streaming vs. Batch and Memory Size.}
We further evaluate the impact of the memory bank size $M$ in Streaming Geo-TTT. With $M=64$, the method has access to only a narrow temporal window and sometimes fails to recover the full manifold structure, leading to a small drop relative to batch mode. Increasing $M$ to $256$ almost closes the gap, while $M=512$ yields no further gains but slightly higher latency. This suggests that Geo-TTT does not require a large replay buffer: a few hundred recent examples are enough to stabilize the manifold across time.



% ==============================================================================
% 7. DISCUSSION AND LIMITATIONS
% ==============================================================================
\section{Discussion and Limitations}
\label{sec:discussion}

While Geo-TTT offers a robust geometric paradigm for adaptation, we acknowledge several limitations that point toward future research directions.

\textbf{Dependency on Batch Size.} 
Unlike standard TTT which can operate on a single stream ($N=1$), Geo-TTT relies on the \textit{collective wisdom} of a batch. If the test batch is too small ($N < 10$) or extremely diverse (no semantic neighbors), the graph becomes sparse, and the Laplacian regularization degrades to an identity mapping (i.e., reverting to Naive TTT). Future work could explore using a \textit{memory bank} to maintain a persistent manifold across streaming data.

\textbf{Computational Overhead.} 
Constructing the $k$-NN graph requires calculating pairwise distances, scaling with $O(N^2)$. For typical LLM inference batch sizes (e.g., $N=64$ or $128$), this cost is negligible ($<1\%$) compared to the Transformer's forward pass. However, for massive-scale batch processing, approximate nearest neighbor (ANN) techniques would be necessary to maintain efficiency.

\textbf{The "Echo Chamber" Risk.} 
Our fundamental assumption is the "Cluster Assumption"—that the majority of neighbors are likely correct. In adversarial scenarios where the majority of the batch is maliciously poisoned or universally biased, Geo-TTT effectively enforces an "echo chamber," potentially suppressing the few correct answers. Robustifying the graph construction against such adversarial attacks remains an open challenge.


% ==============================================================================
% 7. CONCLUSION
% ==============================================================================
\section{Conclusion}
\label{sec:conclusion}

We presented \textbf{Geo-TTT}, a framework that addresses the fragility of Test-Time Training by incorporating geometric inductive biases. By modeling the test batch as a graph on a latent manifold, we showed that \textit{Label Diffusion} acts as a powerful spectral filter, removing high-frequency noise from pseudo-labels. Theoretical bounds and empirical results on LLM reasoning confirm that Geo-TTT effectively breaks the ``Isolation Trap," paving the way for reliable System 2 adaptation.


\section*{Broader Impact}

This work aims to improve the reliability of Large Language Models by making test-time adaptation less fragile and more geometry-aware. By denoising pseudo-labels through manifold-based Label Diffusion, Geo-TTT has the potential to reduce harmful ``hallucinations'' in high-stakes applications such as code generation, scientific assistance, and legal or medical document analysis, where isolated errors can be disproportionately costly. At the same time, our method may also \emph{amplify} existing representational and societal biases: if the latent geometry clusters biased reasoning patterns, enforcing geometric consistency can make those biased predictions more confident and harder to detect. Practitioners should therefore combine Geo-TTT with careful auditing and debiasing procedures, and should not interpret geometric agreement as normative correctness. Finally, like other test-time compute scaling techniques, Geo-TTT increases inference-time energy consumption due to additional gradient steps and graph construction; any deployment should weigh the accuracy gains against the environmental and economic cost of extra compute, and consider lightweight parameter updates or streaming variants to mitigate overhead.

% ==============================================================================
% REFERENCES
% ==============================================================================
\bibliography{refs}
\bibliographystyle{plainnat}

% ==============================================================================
% APPENDIX
% Status: Rigorous Proofs & Details
% ==============================================================================
\clearpage
\appendix
\onecolumn

\section{Detailed Proofs}
\label{app:proofs}

\subsection{Proof of Theorem \ref{thm:denoising} (Geometric Denoising)}

\textbf{Problem Setup.}
Recall the diffusion operator $\mathbf{H} = (\mathbf{I} + \mu \Lap)^{-1}$. Let $\Y_{raw} = \Y_{GT} + \boldsymbol{\xi}$. The refined labels are $\F^* = \mathbf{H} \Y_{raw} = \mathbf{H} \Y_{GT} + \mathbf{H} \boldsymbol{\xi}$.
We analyze the error $\|\F^* - \Y_{GT}\| \le \|\mathbf{H} \Y_{GT} - \Y_{GT}\| + \|\mathbf{H} \boldsymbol{\xi}\|$.

\textbf{Spectral Decomposition.}
Let $\Lap = \sum_{i=1}^N \lambda_i \mathbf{u}_i \mathbf{u}_i^\top$ be the eigendecomposition of the normalized Laplacian, with $0 = \lambda_1 \le \lambda_2 \le \dots \le \lambda_N$.
The operator $\mathbf{H}$ has eigenvalues $h(\lambda_i) = \frac{1}{1 + \mu \lambda_i}$.

\textbf{Term 1: Signal Bias (Smoothing Error).}
By Assumption \ref{ass:spectral}, $\Y_{GT}$ lies in the span of the first $K$ eigenvectors. Thus, $\Y_{GT} = \sum_{i=1}^K c_i \mathbf{u}_i$.
\begin{align}
    \|\mathbf{H} \Y_{GT} - \Y_{GT}\|^2 &= \left\| \sum_{i=1}^K (h(\lambda_i) - 1) c_i \mathbf{u}_i \right\|^2 \\
    &= \sum_{i=1}^K \left( \frac{1}{1+\mu\lambda_i} - 1 \right)^2 c_i^2 = \sum_{i=1}^K \left( \frac{-\mu\lambda_i}{1+\mu\lambda_i} \right)^2 c_i^2
\end{align}
Since $\lambda_i$ are small for $i \le K$ (intra-cluster smoothness), this bias term is small, scaling with $O(\mu^2 \lambda_K^2 \|\Y_{GT}\|^2)$.

\textbf{Term 2: Noise Variance.}
The noise $\boldsymbol{\xi}$ is assumed to be white, i.e., uniformly distributed across the spectrum.
\begin{align}
    \|\mathbf{H} \boldsymbol{\xi}\|^2 &= \sum_{i=1}^N h(\lambda_i)^2 (\mathbf{u}_i^\top \boldsymbol{\xi})^2 \\
    &\le \max_{j} h(\lambda_j)^2 \|\boldsymbol{\xi}\|^2
\end{align}
However, strictly speaking, we gain from the attenuation of high frequencies. If noise is concentrated in high frequencies ($i > K$), then:
\begin{equation}
    \|\mathbf{H} \boldsymbol{\xi}\| \le \frac{1}{1 + \mu \lambda_{K+1}} \|\boldsymbol{\xi}\|
\end{equation}
where $\lambda_{K+1}$ is the first eigenvalue outside the signal subspace (the Spectral Gap). A large gap implies strong noise attenuation.

\textbf{Conclusion.}
Combining the terms:
\begin{equation}
    \|\F^* - \Y_{GT}\| \le \frac{\mu \lambda_K}{1+\mu \lambda_K}\|\Y_{GT}\| + \frac{1}{1+\mu \lambda_{K+1}} \|\boldsymbol{\xi}\|
\end{equation}
This completes the proof. \qed

\section{Additional Experimental Details}
\label{app:details}

\subsection{Hyperparameter Sensitivity}
We perform a grid search for $\mu \in \{0.1, 0.5, 1.0, 2.0, 5.0\}$ and $k \in \{5, 10, 20, 40, 64\}$.
\begin{itemize}
    \item \textbf{Smoothness $\mu$}: We find $\mu=1.0$ is generally robust. Very large $\mu$ ($>5.0$) causes "oversmoothing," where all samples in a batch converge to the same majority class, hurting accuracy.
    \item \textbf{Neighbors $k$}: Performance peaks at $k=20$. For small batch sizes ($N=64$), using $k > 30$ connects semantically unrelated samples, introducing noise into the graph.
\end{itemize}

\subsection{Streaming Implementation}
For Algorithm \ref{alg:streaming}, we use a Memory Bank size of $M=256$. To construct the $k$-NN graph efficiently in streaming mode, we calculate the cosine similarity between the incoming query $z_t$ and all keys in $\mathcal{M}$. We strictly enforce the FIFO policy to handle distribution drift over time.

\section{Additional Ablation: Effect of Memory Size}
\label{app:memory}

In this section, we report the full ablation on the memory bank size $M$ for Streaming Geo-TTT (Algorithm~\ref{alg:streaming}) on ARC-Challenge. We keep all other hyperparameters fixed ($k=20$, $\mu=1.0$) and vary $M$.

\begin{table}[h]
\centering
\caption{Accuracy (\%) on ARC-Challenge for Streaming Geo-TTT with different memory sizes $M$.}
\label{tab:memory}
\vspace{0.1cm}
\begin{tabular}{lccc}
\toprule
Method & $M=64$ & $M=256$ & $M=512$ \\
\midrule
Geo-TTT (Stream) & 54.1 & \textbf{54.8} & 54.9 \\
Geo-TTT (Batch)  & \multicolumn{3}{c}{55.4} \\
\bottomrule
\end{tabular}
\end{table}

We observe that very small memories ($M=64$) under-estimate the test-time manifold, leaving part of the potential gains untapped. Increasing $M$ to $256$ already matches batch-mode Geo-TTT within statistical error, while larger memories do not bring further benefits but incur slightly higher latency. In practice, we therefore recommend setting $M$ in the low hundreds for streaming deployment.


\end{document}